% Results. Every number is a macro from generated/numbers.tex (written by `run.py report` from results/), and every
% table is generated from the same files; no number in this file is typed by hand.

\subsection{Synthetic test set}
Table~\ref{tab:results} compares all systems. On the full test set ($n=\nMethod$ clouds), the method selects the correct family for \accMethod{} \accCIMethod{} of the clouds, with a median Chamfer distance to the noise-free ground-truth shape of \cdMethod{} \cdCIMethod{} (in units of the RMS radius of the cloud). On the baseline subset, where every system sees the same $n=\nBone$ clouds, the method reaches \accMethodsub{}, against \accBtwo{} for the feature classifier B2, \accBone{} for the classical fit-every-family pipeline B1, \accBonefixed{} for its post-hoc repair B1-fixed, and \accBthreea{}--\accBthreec{} for the trivial baselines.

\begin{table}[tbp]
\centering
\caption{Synthetic test set. Fine accuracy is the exact family; coarse accuracy is the shape dimension (point set, curve, surface, blob). Chamfer is the median symmetric Chamfer distance (the sum of the two directed mean distances) between the selected, refined shape and the noise-free ground-truth shape; for surfaces it includes a sampling-grid floor (Sec.~\ref{sec:limitations}). Baselines and the ``baseline subset'' row use the same 20 clouds per family and cell. 95\% percentile cluster-bootstrap CIs over generator chunks; they are conservative (Sec.~\ref{sec:setup}). The full table (superset-tolerant accuracy, residual, axis error) is in \texttt{reports/exp01\_synthetic.md}.}
\label{tab:results}
\resizebox{\textwidth}{!}{\input{generated/results_table_compact.tex}}
\end{table}

\paragraph{Preregistered hypotheses (Table~\ref{tab:hypotheses}).} H1 is supported: the paired macro-F1 difference between the method and B1 is \estHone{} \ciHone{}. H2 is supported: the paired difference in median Chamfer distance is \estHtwo{} \ciHtwo{}. Refinement (H5b) lowers the median Chamfer distance by a small but significant amount, \estHfivebrefinement{} \ciHfivebrefinement{}; because the family is selected by the classifier before refinement, refinement cannot change accuracy (A3 has the same accuracy as the method by construction). The pooled estimate hides opposite effects (post hoc, by the coarse class of the true family): refinement changes the median Chamfer distance by \refineDiffPoint{} \refineCIPoint{} for point sets, \refineDiffCurve{} \refineCICurve{} for curves and \refineDiffSurface{} \refineCISurface{} for surfaces, but by \refineDiffBlob{} \refineCIBlob{} for blobs. A likely reason, not tested: the point-to-sample objective is not an estimate of a Gaussian's covariance. It moves the blob's fixed set of 256 Gaussian draws toward the noisy points, away from the shape the network predicted. H5a (Chamfer versus parameter-space loss) requires the ablation runs, which are postponed (Sec.~\ref{sec:limitations}).

\begin{table}[H]
\centering
\caption{Preregistered hypotheses and post-hoc checks. p-values of H1, H2 and H5b are Holm-corrected over these three synthetic comparisons (H1: normal approximation from the paired cluster-bootstrap standard error; H2, H5b: Wilcoxon signed-rank on per-chunk mean differences); ``$<$ 1e-16'' is a display floor. Rows marked post hoc are outside the Holm family. H3 uses the subsample-majority coarse label, an exact binomial CI and a 10{,}000-permutation null over predicted coarse labels.}
\label{tab:hypotheses}
\input{generated/hypotheses_table.tex}
\end{table}

\paragraph{Where the gap comes from.} Selecting among the network's own refined candidates with geometric AIC (A5) lowers accuracy to \accAfivegaic{}, close to B2. Part of this drop is structural. As implemented, geometric AIC estimates the noise from the best surface fit and therefore can never select the blob (Sec.~\ref{sec:setup}): B1 and A5 recover \blobRecallBone{} and \blobRecallAfivegaic{} of blob clouds, against \blobRecallMethod{} for the method. On clouds whose true family is not blob, A5 still drops to \nonblobAccAfivegaic{} from \nonblobAccMethod{}, and the macro-F1 gap to B1 is \estHoneNonblob{} \ciHoneNonblob{} (post hoc). Against the post-hoc B1-fixed, the macro-F1 difference is \estHoneBonefixed{} \ciHoneBonefixed{} and the median-Chamfer difference \estHtwoBonefixed{} \ciHtwoBonefixed{}. Repairing B1's three handicaps raises its macro-F1 by only \estFoneBonefixedVsBone{} \ciFoneBonefixedVsBone{}. Its accuracy changes by \estAccBonefixedVsBone{} points \ciAccBonefixedVsBone{}, and its median Chamfer distance does not change significantly (Wilcoxon $p=\pCdBonefixedVsBone{}$). B1-fixed recovers \blobRecallBonefixed{} of blob clouds and improves on the largest clouds (\accNlastBonefixed{} vs \accNlastBone{}). Two things offset these gains. Its false blob calls cost recall elsewhere. Allowing as many clusters as points also lets geometric AIC pick degenerate one-point-per-cluster fits on the smallest clouds (\accNfirstBonefixed{} vs \accNfirstBone{}; this explanation came after the run). The method's advantage is therefore not an artefact of B1's implementation defects. Model selection accounts for a large share of the gap. With nested families and realistic noise, a residual-versus-complexity criterion often prefers the more flexible family, while the classifier has learned the prior over families and noise that generated the data. A5 still chooses only among the classifier's top three candidates, so it is not a purely classical selector. Against B1 the accuracy gap grows with the number of points and is smallest for the smallest clouds; against B2 it is largest for the smallest clouds (Fig.~\ref{fig:byn}). The hardest families are cone, helix and paraboloid. The largest single confusions are helix$\leftrightarrow$cone, paraboloid$\leftrightarrow$cone and sphere$\leftrightarrow$cone, but most helix and paraboloid errors go to other families (Appendix, Fig.~\ref{fig:confusion}). Temperature scaling with $T=\temperature$ fitted on the validation stream lowers the expected calibration error from \eceBefore{} to \eceAfter{} and the NLL from \nllBefore{} to \nllAfter{}; the NLL-optimal temperature on the test set itself is \temperatureTest{} (Appendix, Fig.~\ref{fig:reliability}).

\begin{figure}[tbp]
\centering
\includegraphics[width=0.48\textwidth]{fig_by_n.png}\hfill
\includegraphics[width=0.48\textwidth]{fig_by_noise.png}
\caption{Fine accuracy by number of points and by relative noise on the baseline subset (the same clouds for every line). A3 is omitted because its accuracy equals the method's by construction.}
\label{fig:byn}
\end{figure}

\begin{figure}[tbp]
\centering
\includegraphics[width=\textwidth]{fig_gallery_synth.png}
\caption{The first synthetic test cloud of every family in the cell with 65--256 points and relative noise 0.02--0.06: input points (black), ground-truth shape (green), method fit (blue), classical B1 fit (orange).}
\label{fig:gallery}
\end{figure}

\subsection{Real activations}
Table~\ref{tab:real} lists every real dataset and control, and Fig.~\ref{fig:real} shows the clouds with the fitted shapes. As preregistered, H3 uses the majority coarse label over random 80\% subsamples. It agrees with the literature on \hthreeAgree{}/\hthreeN{} preregistered datasets; the label-permutation null gives $p=\hthreeP{}$, so \textbf{H3 is not supported} (it is not falsified either, since agreement is not below 50\%). Post hoc, the full-cloud label would agree on \hthreeAgreeFull{}/\hthreeN{}. Two post-hoc reference points put this in context. Always answering the most frequent expected class (curve) agrees on \hthreeConst{}/\hthreeN{}, so the method does not beat a constant predictor either. At the level of the exact family, the method agrees with the literature on \hthreeFine{}/\hthreeN{} datasets by subsample majority (\hthreeFineFull{}/\hthreeN{} by full-cloud label). The twelve datasets count the toy model R3 twice, with 6 and with 8 features; this split was fixed in the code before extraction and is logged as a deviation.

\paragraph{Coarse agreements.} The datasets in this paragraph agree with the literature in shape dimension; only the first, the 6-feature toy model, also matches the exact family. The toy model of superposition with 6 features (R3) is labelled \fitRthreetmsnsix{} ($p=\pRthreetmsnsix{}$), with a fitted vertex on every feature direction: three orthogonal antipodal pairs (``digons'' in the terminology of \citealp{elhage2022tms}), i.e.\ an octahedron. Llama's character-count manifold (R4) is labelled \fitRfourllamacounting{} ($p=\pRfourllamacounting{}$), a curve as in \citet{gurnee2025linebreaks}; under the preregistered rule a helix claim is not made because ringing is not rejected. Temperature (R8a) is labelled \fitReightallamatemperature{}, the arch PCA produces when similarity decays smoothly with an ordinal variable \citep{novembre2008horseshoe,antognini2018walks}, which agrees with the ``line'' of \citet{bhalla2026sae} at the coarse level; their own 3D PCA view of temperature (their Fig.~2) is also an arch, so ``line'' is a label of their dataset table rather than a description of the projection. Colours (R7) are labelled \fitRsevenllamacolors{}, a curved surface as reported (paraboloid in \citealp{bhalla2026sae}), although the cloud carries the $\lambda_3/\lambda_4<1.5$ flag, so its third principal axis is not well separated from the fourth.

\paragraph{Misses and their likely causes (post-hoc inspection, not tested).} At the coarse level, H3 is missed on five datasets. (i) GPT-2's positional embeddings (R1) are a 1D spiral whose radius varies along the curve; no family represents it (helix has constant radius, ringing three harmonics), and the classifier selects \fitRonegpttwowpos{} ($p=\pRonegpttwowpos{}$), a surface the spiral lies on. The held-out residual cannot reveal this error because held-out points also lie on the sphere. (ii) Years (R6) keep only \varRsixllamayears{} of the variance in the top three components and are flagged ($\lambda_3/\lambda_4<1.5$): the 3D view cuts through a higher-dimensional structure. The 3D PCA view of years in \citet{bhalla2026sae} (their Fig.~2) is itself a broad cloud; ``helix'' is the label in their dataset table, which the text does not describe further. (iii) Age (R8b) is not one smooth curve: ordered by age, consecutive points jump apart around round decade values; it is labelled \fitReightbllamaage{}. (iv) True/false city statements (R9) form two overlapping clouds of unequal spread, outside the vocabulary (labelled \fitRninellamatruthcities{}). (v) The 8-feature toy model (R3) is labelled \fitRthreetmsneight{} on the full cloud ($p=\pRthreetmsneight{}$), but most 80\% subsamples, which keep 6 of its 8 vertices, are labelled \majRthreetmsneight{} (\majfreqRthreetmsneight{}), so it counts as a miss under the preregistered rule. At the fine level, two more failures do not affect H3: weekdays and months (R5) form one tight cluster per token, 7 and 12 clusters, more than the 6 centroids the clusters family allows, and the fitter threads a \fitRfiveallamaweekdays{} through them; the grokked embeddings (R2) keep only \varRtwogrokkingWE{} of the variance in the top three components, as expected for the several Fourier circles of a grokked model \citep{nanda2023progress}, and are labelled \fitRtwogrokkingWE{} instead of a circle.

\begin{figure}[tbp]
\centering
\includegraphics[width=\textwidth]{fig_real.png}
\caption{Real activation clouds in their top-3 PCA coordinates, coloured by concept value, with the fitted shape (red; crosses for point-set vertices). Titles give the full-cloud fit with its calibrated probability and the majority label over 200 random 80\% subsamples.}
\label{fig:real}
\end{figure}

\paragraph{Controls (H4).} Matched-covariance Gaussians are labelled blob in \hfourBlob{} of draws (mean over the \hfourNreal{} real datasets, exploratory ones included; per dataset the rate falls below the preregistered 80\% for the smallest clouds, the toy models and per-concept means, see \texttt{reports/deviations.md}), and the ordinal-process controls ringing in \hfourRing{} versus helix in \hfourHelix{}, so H4 passes under the implemented criterion (the mean blob rate rather than a per-dataset 80\% threshold). The preregistered wording (``C2 is labelled blob in $\ge$ 80\% of subsamples'') admits both readings. Pooled over the preregistered datasets the rate is \hfourBlobPrereg{}, which passes. Read per dataset, only \hfourPreregOk{} of the \hfourPreregN{} preregistered datasets reach 80\%, so H4 fails under that reading. A post-hoc split shows that this pass relies on the random walks (ringing in \hfourWalkRing{}): the smooth Gaussian-process controls alone are labelled helix in \hfourGpHelix{}. The classifier therefore confuses smooth ordinal processes with helices, the confound the preregistered helix rule guards against; for the fully evaluated Gaussian-process control the held-out residual favours ringing and the rule rejects the helix claim. Post hoc, the same rule was applied to every full-cloud helix fit with held-out folds: it rejects the helix claim for \helixRuleRejected{} but accepts it for \helixRuleAccepted{}, all of which the literature describes as circles. The rule separates helices from ringing, not from circles or from clusters threaded by a helix. Randomly initialised GPT-2 positional embeddings and grokking embeddings are labelled blob, but randomly initialised Llama is labelled \fitCthreellamamonthsinit{} for the month prompts and \fitCthreellamaweekdaysinit{} for the weekday prompts: Fig.~\ref{fig:controls} shows genuine clusters, one per input token, because a random transformer still separates token identities. The preregistered expectation ``blob'' for this control was wrong, not the fitter.

\begin{figure}[tbp]
\centering
\includegraphics[width=\textwidth]{fig_controls.png}
\caption{Controls: a smooth Gaussian process and a random walk over an ordinal index (C1), and network controls (C3): randomly initialised GPT-2, Llama and grokking weights, the grokking model at memorisation, and a toy model trained on dense features.}
\label{fig:controls}
\end{figure}

\begin{sidewaystable}
\centering
\caption{Real datasets and controls (C1 shows the first of 10 draws of each process). Flags: $\lambda_3/\lambda_4<1.5$ (gap); OOD: residual above the synthetic 95th percentile of correct fits at that $N$. Held-out residuals, matched-Gaussian rates and model/layer details are in the full table in \texttt{reports/exp02\_real\_activations.md}.}
\label{tab:real}
\input{generated/real_table_compact.tex}
\end{sidewaystable}
